\documentclass[t,usepdftitle=false,aspectratio=169]{beamer}

\usepackage[utf8]{inputenc}
\usetheme{Singapore}
\usepackage{xcolor}
\setbeamertemplate{footline}[frame number]

% \setbeamercovered{transparent}
%\usecolortheme{crane}
\title[IFT3515]{IFT 3515\\Fonctions à plusieurs variables\\Optimisation avec contraintes\\Méthodes de pénalité}
\author[Fabian Bastin]{Fabian Bastin\\DIRO\\Université de Montréal}
\date{}

\usepackage{enumerate}
\usepackage[francais]{babel}

\usepackage{easybmat}
\usepackage{graphicx}
\usepackage{tikz}
\usetikzlibrary{arrows.meta,calc,decorations.markings}

\newtheorem{defn}{Définition}
\newtheorem{lem}{Lemme}
\newtheorem{thm}{Théorème}
\newtheorem{coro}{Corollaire}

\def\red{\color{red}}
\def\blue{\color{blue}}

\def\co{\mathcal{o}}

\def\cA{\mathcal{A}}
\def\cB{\mathcal{B}}
\def\cE{\mathcal{E}}
\def\cI{\mathcal{I}}
\def\cL{\mathcal{L}}
\def\cN{\mathcal{N}}
\def\cT{\mathcal{T}}
\def\cR{\mathcal{R}}
\def\cX{\mathcal{X}}
\def\cV{\mathcal{V}}

\def\bu{\boldsymbol{u}}

\def\NN{\mathbb{N}}
\def\RR{\mathbb{R}}

\begin{document}
\frame{\titlepage}

% ------------------------------------------------------------------------------------------------------------------------------------------------------

\begin{frame}
\frametitle{Motivations}

%Partiellement basé sur \url{https://courses.maths.ox.ac.uk/node/view_material/1377}

\mbox{}

De nombreux problèmes présentent un ensemble réalisable sur lequel il est difficile de calculer des projections, même si cet ensemble est convexe.

\mbox{}

Nous allons tenter de construire des problèmes plus simples, mais permettant d'approcher les solutions des véritables problèmes avec une précision arbitraire (au moins en théorie).

\end{frame}

\begin{frame}
\frametitle{Problèmes non-linéaires avec contraintes d'égalité}

Nous nous intéressons dans un premier temps aux problèmes présentant uniquement des contraintes d'égalité, de la forme
\begin{align*}
\min_{x \in \RR^n}\ & f(x) \\
\mbox{t.q. } &g_i(x) = 0,\ i \in \cE
\end{align*}
où $f: \RR^n \rightarrow \RR$, $g_i: \RR^n \rightarrow \RR$, $i \in \cE$. De plus, nous supposons que $f$ et $g_i$ ($i \in \cE$) sont des fonction ``douces'' (typiquement deux fois continûment différentiables).

\mbox{}

Notre but est de trouver des solutions locales, ou à tout le moins des points KKT.

\mbox{}

Dans le cas de contraintes d'égalité générales (non linéaires), il est souvent très difficile, voire impossible, de générer des points réalisables.

\end{frame}

\begin{frame}
\frametitle{Problèmes non-linéaires avec contraintes d'égalité}

Nous allons déplacer les contraintes dans l'objectif et former un nouveau problème sans contraintes, paramétrisé pour tenir compte du double objectif de minimisation et de satisfaction des contraintes.

\mbox{}

Regroupons les contraintes en une multifonction $g: \RR^n \rightarrow \RR^m$, avec $m = \#\cE$:
$$
g(x) = \begin{pmatrix} g_1(x) \\ g_2(x) \\ \vdots \\ g_m(x) \end{pmatrix}
$$
Le problème d'optimisation se réécrit
\begin{align*}
\min_{x \in \RR^n}\ & f(x) \\
\mbox{t.q. } &g(x) = 0.
\end{align*}

\end{frame}

\begin{frame}
\frametitle{Pénalisation quadratique}

Le problème avec pénalité quadratique s'écrit
$$
\min_{x \in \RR^n} \Phi_{\mu}(x) = f(x) + \frac{1}{2\mu} \| g(x) \|^2
$$
où $\mu > 0$ est le paramètre de pénalisation.

\mbox{}

Nous pouvons voir le nouveau problème comme un problème bi-objectif où $\mu$ pénalise la violation des contraintes, et $\mu \rightarrow 0$ force les contraintes à être satisfaites, tout en atteignant l'optimalité sur $f$.

\mbox{}

$\Phi_{\mu}$ peut cependant avoir d'autres points stationnaires qui ne sont pas solutions du problème initial, par exemple si $g(x) = 0$ est non consistant (il n'y a pas de solution réalisable pour le problème initial).

\end{frame}

\begin{frame}
\frametitle{Méthode de pénalisation quadratique}

Étant donné $\mu^0 > 0$, soit $k = 0$.
Jusqu'à ``convergence'', répéter
\begin{enumerate}
\item
Choisir $0 < \mu^{k+1} < \mu^k$.
\item
En partant de $x_0^k$ (possiblement $x_0^k := x^k$), utiliser un algorithme de minimisation sans contrainte pour trouver un minimiseur approximatif $x^{k+1}$ de $\Phi_{\mu^{k+1}}$.
Poser $k := k+1$.
\end{enumerate}

On doit avoir $\mu^k \rightarrow 0$ comme $k \rightarrow \infty$.
Par exemple: $\mu^{k+1} := 0.1\mu^k$, $\mu^{k+1} := (\mu^k)^2$, etc.

\mbox{}

Algorithmes de minimisation
\begin{itemize}
\item
recherche linéaire, méthode de région de confiance
\item
$\mu$ petit: $\Phi_{\mu}$ est très raide dans la direction des contraintes des gradients, aussi observons-nous de rapides changements de $\Phi_{\mu}$ dans ces directions; implication pour la ``forme'' de la région de confiance.
\end{itemize}

\end{frame}

\begin{frame}
\frametitle{Convergence}

\begin{thm}[Convergence globale de la méthode de pénalité]
Appliquons la méthode de pénalisation quadratique au programme avec contraintes d'égalité.
Supposons que $f$, $g \in C^1$,
$$
y_i^k = \frac{g_i(x^k)}{\mu^k}, i \in \cE,
$$
et
$$
\left\| \nabla \Phi_{\mu^k}(x^k) \right\| \leq \epsilon^k,
$$
où $\epsilon^k \rightarrow 0$, comme $k \rightarrow \infty$
et $\mu^k \rightarrow 0$, comme $k \rightarrow \infty$.

Supposons de plus que $x^k \rightarrow x^*$, où $\nabla g_i(x^*)$, $i \in \cE$, sont linéairement indépendants.

Alors $x^*$ est un point KKT du problème avec contraintes d'égalité, et $y^k \rightarrow y^*$, où $y^*$ est le vecteur de multiplicateurs de Lagrange des contraintes d'égalité.
\end{thm}

\end{frame}

\begin{frame}
\frametitle{Dérivées de la fonction de pénalité}

\begin{itemize}
\item
Soit
$$
y(\mu) := \frac{g(x)}{\mu}
$$
des estimateurs des multiplicateurs de Lagrange.
\item
Soit $L$ la fonction lagrangienne du problème avec contraintes d'égalité:
$$
L(x,y) := f(x) + y^Tg(x)
$$
\item
Le gradient de la fonction pénalisée
$$
\Phi_{\mu}(x) = f(x) + \frac{1}{2\mu} \| g(x) \|^2
$$
vaut
$$
\nabla \Phi_{\mu}(x) = \nabla f(x) + \frac{1}{\mu} J(x)^T g(x) = \nabla_x L(x, y(\mu)),
$$
où $J(x)$ est la matrice jacobienne, de dimensions $m \times n$, des contraintes $g(x)$.
\end{itemize}

\end{frame}

\begin{frame}
\frametitle{Dérivées de la fonction de pénalité}

\begin{itemize}
\item
La matrice hessienne de $\Phi_{\mu}$ est
\begin{align*}
\nabla^2 \Phi_{\mu}(x)
&= \nabla^2 f(x) + \frac{1}{\mu} \sum_{i \in \cE} g_i(x) \nabla^2 g_i(x) + \frac{1}{\mu} J(x)^TJ(x) \\
&= \nabla^2_{xx} L(x, y(\mu)) + \frac{1}{\mu} J(x)^TJ(x)
\end{align*}
\item
Comme $\mu \rightarrow 0$,
\begin{itemize}
\item
généralement, $g_i(x) \rightarrow 0$ ou $\nabla^2 g_i(x) \rightarrow 0$ au même taux avec $\mu$ pour tous les indices $i$.
Dès lors, généralement, $\nabla^2_{xx} L(x, y(\mu))$ se comporte bien.
\item
Mais
$$
\frac{1}{\mu} J(x)^TJ(x) \rightarrow \frac{1}{0} J(x^*)^TJ(x^*) = \infty
$$
\end{itemize}
\end{itemize}

\end{frame}

% ---------- frame 1a : référence aux cônes, spécialisation ----------
\begin{frame}
\frametitle{Mauvais conditionnement: cônes tangent et normal}

\begin{columns}[T]
\begin{column}{0.55\textwidth}
Le jeu de transparents sur les méthodes de projection introduit les cônes $T_{\cX}(x^*)$ et $N_{\cX}(x^*)$ dans leur forme générale. Le jeu sur les conditions d'optimalité établit ensuite que sous la LICQ:
$$
T_{\cX}(x^*) = \ker J(x^*),\quad N_{\cX}(x^*) = \mathrm{Im}\,J(x^*)^T.
$$
$J^TJ\,v = 0 \Leftrightarrow Jv = 0$, donc $\ker(J^TJ) = \ker J$:
\begin{itemize}
  \item val.\ prop.\ \textbf{nulles}: $v \in T_{\cX}(x^*)$ \textbf{(tangentes)}
  \item val.\ prop.\ $\sigma_i^2>0$: $v \in N_{\cX}(x^*)$ \textbf{(normales)}
\end{itemize}
\end{column}
\begin{column}{0.42\textwidth}
\centering
\begin{tikzpicture}[scale=1.5,>=Stealth]
  % Hyperbola x1*x2=1, x1 in [0.35,2.6]
  \draw[thick,blue!70!black,domain=0.38:2.55,samples=80,smooth]
    plot (\x, {1/\x});
  \node[blue!70!black,right,font=\small] at (2.1,0.52) {$x_1x_2=1$};
  % Axes
  \draw[->] (0,0) -- (2.7,0) node[right,font=\small] {$x_1$};
  \draw[->] (0,0) -- (0,2.7) node[above,font=\small] {$x_2$};
  % x* = (1,1)
  \fill[red] (1,1) circle (1.8pt) node[below right,font=\small] {$x^*$};
  % Tangent direction: (1,-1)/sqrt(2), scaled by 0.75
  \draw[->,thick,green!60!black] (1,1) -- ++(0.53,-0.53)
    node[right,font=\small,green!60!black] {$T_{\cX}$};
  \draw[->,thick,green!60!black] (1,1) -- ++(-0.53,0.53);
  % Normal direction: (1,1)/sqrt(2), scaled by 0.65
  \draw[->,thick,orange!80!black] (1,1) -- ++(0.46,0.46)
    node[right,font=\small,orange!80!black] {$N_{\cX}$};
  % Tangent line (dashed)
  \draw[dashed,green!50!black,thin] (0.35,1.65) -- (1.65,0.35);
\end{tikzpicture}
\end{column}
\end{columns}

\end{frame}

\begin{frame}
\frametitle{Mauvais conditionnement: analyse spectrale}

Dans la décomposition $\nabla^2\Phi_\mu = \nabla^2_{xx}L(x,y(\mu)) + \frac{1}{\mu}J^TJ$, pour $\mu \rightarrow 0$:
\begin{itemize}
  \item directions \textbf{normales} $v \in N_{\cX}(x^*)$ ($Jv \ne 0$):
  $$\frac{1}{\mu}J^TJ\,v = \frac{\sigma^2}{\mu}\,v \;\longrightarrow\; +\infty$$
  La pénalité pénalise de plus en plus toute sortie de $\cX$ dans ces directions.
  \item directions \textbf{tangentes} $v \in T_{\cX}(x^*)$ ($Jv = 0$):
  $$\frac{1}{\mu}J^TJ\,v = 0, \quad \text{contribution } O(1) \text{ de } \nabla^2_{xx}L$$
  Ces directions restent dans $\cX$ au premier ordre: la pénalité ne les affecte pas.
\end{itemize}

\mbox{}

Les $m = \dim N_{\cX}(x^*)$ valeurs propres divergent en $O(1/\mu)$; les $n-m = \dim T_{\cX}(x^*)$ restantes sont en $O(1)$. Le nombre de conditionnement est donc en $O(1/\mu) \rightarrow \infty$.

\end{frame}

% ---------- frame 2/3 : exemple numérique ----------
\begin{frame}
\frametitle{Mauvais conditionnement: exemple numérique}

\begin{columns}[T]
\begin{column}{0.52\textwidth}
$\min \tfrac{1}{2}(x_1^2+x_2^2)$ s.c.\ $x_1x_2=1$, $x^*\!=\!(1,1)^T$.

En $x^*$: $\nabla^2\Phi_\mu = \begin{pmatrix}1+\tfrac{1}{\mu}&\tfrac{1}{\mu}\\[1pt]\tfrac{1}{\mu}&1+\tfrac{1}{\mu}\end{pmatrix}.$

$\lambda_1=1+\tfrac{2}{\mu}$ ($\in N_{\cX}$), $\lambda_2=1$ ($\in T_{\cX}$).

\begin{center}\small
\begin{tabular}{c|c|c}
$\mu$ & $\lambda_1$ & $\kappa$ \\\hline
$1$     & $3$    & $3$    \\
$0.3$   & $\approx 8$ & $\approx 8$ \\
$0.1$   & $21$   & $21$   \\
\end{tabular}
\end{center}
$\kappa \sim 2/\mu \rightarrow \infty$: Newton s'écrase dans $N_{\cX}$.
\end{column}
\begin{column}{0.45\textwidth}
\centering
% Ellipses of level set d^T H d = 1
% H eigenvectors: v1=(1,1)/sqrt2 (normal), v2=(1,-1)/sqrt2 (tangent)
% semi-axes: a=1/sqrt(lambda1) along v1, b=1/sqrt(lambda2)=1 along v2
% Rotation by 45 deg: axes along (1,1) and (1,-1)
\begin{tikzpicture}[scale=1.1]
  % Axes with labels for directions
  \draw[->] (-1.3,0)--(1.3,0) node[right,font=\tiny] {};
  \draw[->] (0,-1.3)--(0,1.3) node[above,font=\tiny] {};
  % Diagonal axis labels
  \draw[->,gray!60,thin] (-1.1,-1.1)--(1.1,1.1)
    node[above right,font=\tiny,gray!60] {$N_{\cX}$};
  \draw[->,gray!60,thin] (1.1,-1.1)--(-1.1,1.1)
    node[above left,font=\tiny,gray!60] {$T_{\cX}$};
  % mu=1: a=0.577, b=1 rotated 45deg
  % ellipse in rotated frame: (u/a)^2+(v/b)^2=1, u along (1,1)/sqrt2, v along (1,-1)/sqrt2
  % parametric: x=a*cos(t)*(1,1)/sqrt2 + b*sin(t)*(1,-1)/sqrt2
  \draw[blue,thick,domain=0:360,samples=120,smooth]
    plot ({0.577*cos(\x)*0.7071 + 1.0*sin(\x)*0.7071},
         {0.577*cos(\x)*0.7071 - 1.0*sin(\x)*0.7071})
    node[above right,font=\tiny,blue] {$\mu\!=\!1$};
  % mu=0.3: a=0.361, b=1
  \draw[orange,thick,domain=0:360,samples=120,smooth]
    plot ({0.361*cos(\x)*0.7071 + 1.0*sin(\x)*0.7071},
         {0.361*cos(\x)*0.7071 - 1.0*sin(\x)*0.7071})
    node[right,font=\tiny,orange] {$\mu\!=\!0.3$};
  % mu=0.1: a=0.218, b=1
  \draw[red,thick,domain=0:360,samples=120,smooth]
    plot ({0.218*cos(\x)*0.7071 + 1.0*sin(\x)*0.7071},
         {0.218*cos(\x)*0.7071 - 1.0*sin(\x)*0.7071})
    node[right,font=\tiny,red] {$\mu\!=\!0.1$};
  % Origin
  \fill (0,0) circle (1.5pt);
\end{tikzpicture}

{\small\centering Courbes de niveau de $d^T\nabla^2\Phi_\mu\, d=1$}
\end{column}
\end{columns}

\end{frame}

% ---------- frame 3/3 : conséquence pour Newton ----------
\begin{frame}
\frametitle{Mauvais conditionnement: conséquence pour Newton}

Il n'est pas clair que nous puissions calculer précisément les changements sur $x^k$.
En particulier, asymptotiquement, nous souhaitons minimiser $\Phi_{\mu^{k+1}}(x)$ au moyen d'un pas de Newton, i.e. en résolvant
$$
\nabla^2 \Phi_{\mu}(x)\, s = -\nabla \Phi_{\mu}(x)
$$
pour obtenir le pas $s$ à partir de la solution courante $x = x_i^k$, $\mu = \mu^{k+1}$.

\mbox{}

Pour l'exemple précédent ($x_1x_2=1$), à $x^* = (1,1)^T$ et $\mu = 0.001$, $\kappa \approx 2001$: les erreurs d'arrondi en flottant double ($\epsilon_{\rm mach} \approx 10^{-16}$) se propagent en $\kappa\,\epsilon_{\rm mach} \approx 2\times 10^{-13}$.

\mbox{}

Malgré ce mauvais conditionnement apparent, il est toujours possible de calculer $s$ avec précision en reformulant le système (voir transparent suivant).

\end{frame}

\begin{frame}
\frametitle{Calcul du pas de Newton}

En utilisant les calculs de dérivées, le système
$$
\nabla^2 \Phi_{\mu}(x) s = -\nabla \Phi(x)
$$
est équivalent à
$$
\left( \nabla^2_{xx} L(x, y(\mu)) + \frac{1}{\mu} J(x)^TJ(x) \right) s = - \left( \nabla f(x) + \frac{1}{\mu} J(x)^T g(x) \right)
$$
où $y(\mu) = g(x)/\mu$.
Définissons la variable auxiliare $w$
$$
w = \frac{1}{\mu} \left( J(x)s + g(x) \right)
$$
Alors, le système de Newton peut être réécrit comme
$$
\begin{pmatrix}
\nabla^2_{xx} L(x, y(\mu)) & J(x)^T \\
J(x) & -\mu I
\end{pmatrix}
\begin{pmatrix}
s \\ w
\end{pmatrix}
=
- \begin{pmatrix}
\nabla f(x) \\ g(x)
\end{pmatrix}
$$

\end{frame}

\begin{frame}
\frametitle{Calcul du pas de Newton}

Ce système est essentiellement indépendant de $\mu$ pour de petites valeurs de $\mu$, et dès lors ne souffre pas de mauvais conditionnement comme $\mu \rightarrow 0$.

\mbox{}

Il faut toujours être prudent lors de la minimisation de
$\Phi_{\mu}$ pour de petits $\mu$.

\mbox{}

Par exemple, en utilisant des régions de confiance, il est bon d'utiliser comme test de contrainte de région de confiance $\| s \|_B \leq \Delta$, où $B$ prend en compte les termes mal-conditionnés de la Hessienne afin d'encourager une décroissance égale du modèle dans toutes les directions.

\end{frame}

\begin{frame}
\frametitle{Conditions d'optimalité perturbées}

Reprenons le problème
\begin{align*}
\min_{x \in \RR^n}\ & f(x) \\
\mbox{t.q. } &g_i(x) = 0,\ i \in \cE
\end{align*}

Des conditions KKT, nous pouvons considérer le système
$$
\begin{cases}
\nabla f(x) + J(x)^Ty = 0 & \mbox{(annulation du gradient du Lagrangien)} \\
g(x) = 0 & \mbox{(faisabilité primale)}
\end{cases}
$$
Considérons le problème perturbé
$$
\begin{cases}
\nabla f(x) + J(x)^Ty = 0 & \mbox{(annulation du gradient du Lagrangien)} \\
g(x) - \mu y = 0 & \mbox{(faisabilité primale)}
\end{cases}
$$
Pour trouver les racines de ce système non-linéaire comme $\mu \rightarrow 0$ ($\mu > 0$), nous pouvons utiliser la méthode de Newton de recherche de racine.

\end{frame}

\begin{frame}
\frametitle{Conditions d'optimalité perturbées}

La méthode de Newton pour le problème pertubé calcule les changements $(\Delta x, \Delta y)$ en résolvant le système
$$
\begin{pmatrix}
\nabla^2_{xx} L(x, y) & J(x)^T \\
J(x) & -\mu I
\end{pmatrix}
\begin{pmatrix}
\Delta x \\ \Delta y
\end{pmatrix}
=
- \begin{pmatrix}
\nabla f(x) + J(x)^Ty \\ g(x) - \mu y
\end{pmatrix}
$$
donnant
\begin{align*}
\nabla^2_{xx} L(x, y)\Delta x + J(x)^T\Delta y &=  -\nabla f(x) - J(x)^Ty \\
J(x)\Delta x -\mu \Delta y &= -g(x) + \mu y
\end{align*}
La deuxième équation permet d'obtenir
$$
\Delta y = \frac{1}{\mu}J(x)\Delta x + \frac{1}{\mu}g(x) - y
$$

\end{frame}

\begin{frame}
\frametitle{Conditions d'optimalité perturbées}

En remplaçant $\Delta y$ dans l'égalité précédente, nous avons
\begin{align*}
&\nabla^2_{xx} L(x, y)\Delta x + J(x)^T\left( \frac{1}{\mu}J(x)\Delta x + \frac{1}{\mu}g(x) - y \right) \\ &=  -\nabla f(x) - J(x)^Ty
\end{align*}
ce qui donne
$$
\left( \nabla^2_{xx} L(x, y) + \frac{1}{\mu} J(x)^TJ(x) \right) \Delta x = - \left( \nabla f(x) + \frac{1}{\mu} J(x)^T g(x) \right)
$$
Similaire au système de Newton pour la pénalité quadratique!

\end{frame}

\begin{frame}
\frametitle{Conditions d'optimalité perturbées}

{\blue Primal:}
$$
\left( \nabla^2_{xx} L(x, y(\mu)) + \frac{1}{\mu} J(x)^TJ(x) \right) s = - \left( \nabla f(x) + \frac{1}{\mu} J(x)^T g(x) \right)
$$
où $y(\mu) = g(x)/\mu$.

\mbox{}

{\blue Primal-Dual:}
$$
\left( \nabla^2_{xx} L(x, y) + \frac{1}{\mu} J(x)^TJ(x) \right) \Delta x = - \left( \nabla f(x) + \frac{1}{\mu} J(x)^T g(x) \right)
$$

\mbox{}

La différence réside dans la liberté de choisir $y$ dans $\nabla^2_{xx} L(x, y)$ dans les méthodes primales-duales.
En termes computationnels, cela fait une différence importante.

\end{frame}

\begin{frame}
\frametitle{Intérêt de l'approche primale-duale: mise en place}

Reprenons: $\min \tfrac{1}{2}(x_1^2+x_2^2)$ s.c.\ $g(x)=x_1x_2-1=0$.

Le Lagrangien est $L(x,y) = \tfrac{1}{2}(x_1^2+x_2^2) + y(x_1x_2-1)$, d'où:
$$
\nabla^2 f = I_2, \quad \nabla^2 g = \begin{pmatrix}0&1\\1&0\end{pmatrix},
\quad
\nabla^2_{xx}L(x,y) = I_2 + y\begin{pmatrix}0&1\\1&0\end{pmatrix} = \begin{pmatrix}1&y\\y&1\end{pmatrix}.
$$

Les valeurs propres de $\begin{pmatrix}1&y\\y&1\end{pmatrix}$ vérifient $(1-\lambda)^2-y^2=0$, donc $\lambda = 1\pm|y|$. La matrice est définie positive ssi $|y|<1$, singulière si $|y|=1$, indéfinie si $|y|>1$.

\mbox{}

Les conditions KKT en $x^*$ donnent $\nabla f(x^*) + y^*\nabla g(x^*)=0$:
$$
\begin{pmatrix}1\\1\end{pmatrix} + y^*\begin{pmatrix}x_2^*\\x_1^*\end{pmatrix}
= \begin{pmatrix}1\\1\end{pmatrix} + y^*\begin{pmatrix}1\\1\end{pmatrix} = 0
\implies y^* = -1,
$$
ce qui correspond bien à $|y^*|=1$: la Hessienne du Lagrangien est singulière à la solution (cas limite).

\end{frame}

\begin{frame}[shrink=2]
\frametitle{Intérêt de l'approche primale-duale: approche primale}

Itéré infaisable: $x^k=(2,1)^T$, $g(x^k) = 2\cdot 1-1 = 1$.

Dans l'approche \textbf{primale}, le multiplicateur est entièrement déterminé par la pénalisation:
$$
y(\mu) = \frac{g(x^k)}{\mu} = \frac{1}{\mu} \xrightarrow{\mu\to 0} +\infty.
$$

Ainsi $\nabla^2_{xx}L(x^k,y(\mu)) = \begin{pmatrix}1&1/\mu\\1/\mu&1\end{pmatrix}$, de valeurs propres $1\pm 1/\mu$:
\begin{center}
\begin{tabular}{c|c|c|c}
$\mu$ & $y(\mu)=1/\mu$ & val.\ prop.\ $1\pm y(\mu)$ & statut de $\nabla^2_{xx}L$ \\\hline
$1$    & $1$   & $0,\ 2$       & singulière \\
$0.5$  & $2$   & $-1,\ 3$     & \textcolor{red}{indéfinie} \\
$0.1$  & $10$  & $-9,\ 11$    & \textcolor{red}{très indéfinie} \\
$0.01$ & $100$ & $-99,\ 101$  & \textcolor{red}{diverge} \\
\end{tabular}
\end{center}

À partir de $\mu < 1$ (i.e.\ $|y(\mu)|>1$), la Hessienne du Lagrangien devient indéfinie, ce qui compromet la résolution du système de Newton et la descente du modèle quadratique.

\end{frame}

\begin{frame}
\frametitle{Approche primale-duale: itération de Newton}

Le système KKT est résolu par Newton. À chaque itération, on résout:
$$
\begin{pmatrix}\nabla^2_{xx}L(x^k,y^k) & J(x^k)^T \\ J(x^k) & 0\end{pmatrix}
\begin{pmatrix}\Delta x \\ \Delta y\end{pmatrix}
= -\begin{pmatrix}\nabla_x L(x^k,y^k) \\ g(x^k)\end{pmatrix}
$$
puis $x^{k+1} = x^k + \Delta x$, $y^{k+1} = y^k + \Delta y$.

\mbox{}

\textbf{Itération $k=0$:} $x^0=(2,1)^T$, $y^0=0$.

$J(x^0) = (1,\,2)$, $\nabla^2_{xx}L(x^0,0) = I_2$, $g(x^0)=1$, $\nabla_x L(x^0,0) = (2,1)^T$.

Le système de Newton $K^0(\Delta x, \Delta y)^T = -(2,1,1)^T$ s'écrit:
$$
\begin{pmatrix}1&0&1\\0&1&2\\1&2&0\end{pmatrix}\begin{pmatrix}\Delta x_1\\\Delta x_2\\\Delta y\end{pmatrix} = \begin{pmatrix}-2\\-1\\-1\end{pmatrix}.
$$

\end{frame}

\begin{frame}
\frametitle{Approche primale-duale: itération $k=0$ (suite)}

Des lignes 1 et 2: $\Delta x_1 = -2-\Delta y$, $\Delta x_2 = -1-2\Delta y$. En substituant dans la ligne 3:
$$
(-2-\Delta y)+2(-1-2\Delta y) = -1 \;\Longrightarrow\; -5\Delta y = 3 \;\Longrightarrow\; \Delta y = -\tfrac{3}{5}.
$$
Donc $\Delta x_1=-\tfrac{7}{5}$, $\Delta x_2=\tfrac{1}{5}$, et
$$
x^1=\bigl(\tfrac{3}{5},\,\tfrac{6}{5}\bigr)^T = (0.6,\,1.2)^T, \qquad y^1=-\tfrac{3}{5} = -0.6.
$$

On vérifie: $g(x^1)=\tfrac{3}{5}\cdot\tfrac{6}{5}-1 = \tfrac{18}{25}-1 = -\tfrac{7}{25}$, $\nabla^2_{xx}L(x^1,y^1)=\begin{pmatrix}1&-0.6\\-0.6&1\end{pmatrix}$ ($\kappa=4$, \textcolor{blue}{$\succ 0$}).

\end{frame}

\begin{frame}
\frametitle{Approche primale-duale: itération de Newton (suite)}

\textbf{Itération $k=1$:} $x^1=(\tfrac{3}{5},\tfrac{6}{5})^T$, $y^1=-\tfrac{3}{5}$.
$$
J(x^1)=\bigl(\tfrac{6}{5},\tfrac{3}{5}\bigr),\quad
\nabla^2_{xx}L = \begin{pmatrix}1&-\tfrac{3}{5}\\-\tfrac{3}{5}&1\end{pmatrix},\quad
g(x^1)=\tfrac{18}{25}-1=-\tfrac{7}{25}.
$$
$$
\nabla_x L = \begin{pmatrix}\tfrac{3}{5}\\[2pt]\tfrac{6}{5}\end{pmatrix}
+ \Bigl(-\tfrac{3}{5}\Bigr)\begin{pmatrix}\tfrac{6}{5}\\[2pt]\tfrac{3}{5}\end{pmatrix}
= \begin{pmatrix}-\tfrac{3}{25}\\[2pt]\tfrac{21}{25}\end{pmatrix}.
$$

Résolution numérique: $\Delta x \approx (0.407, -0.348)^T$, $\Delta y \approx -0.413$.
$$
x^2 \approx (1.007,\; 0.852)^T, \qquad y^2 \approx -1.013.
$$

\mbox{}

Après $k=1$, $x^2$ est quasi-faisable ($g(x^2)\approx -0.142$) et $y^2 \approx y^*$: la convergence quadratique s'amorce.

\end{frame}

\begin{frame}
\frametitle{Approche primale-duale: convergence}

\begin{columns}[T]
\begin{column}{0.50\textwidth}
\begin{center}\small
\begin{tabular}{c|c|c|c}
$k$ & $x_1^k$ & $x_2^k$ & $g(x^k)$ \\\hline
0 & 2.0000 & 1.0000 &  $+1.000$ \\
1 & 0.6000 & 1.2000 &  $-0.280$ \\
2 & 1.0072 & 0.8523 &  $-0.142$ \\
3 & 1.0000 & 0.9989 &  $-0.001$ \\
4 & 1.0000 & 1.0000 &  $\approx 0$ \\
\end{tabular}
\end{center}
\vspace{2pt}
\begin{itemize}\small
  \item 4 itérations depuis un point infaisable.
  \item Convergence quadratique après $k=2$.
  \item $y^k \to y^*=-1$ indépendamment de $\mu$.
\end{itemize}
\end{column}
\begin{column}{0.47\textwidth}
\centering
\begin{tikzpicture}[scale=1.35,>=Stealth]
  % Axes
  \draw[->] (0,0)--(2.5,0) node[right,font=\small] {$x_1$};
  \draw[->] (0,0)--(0,2.0) node[above,font=\small] {$x_2$};
  % Hyperbola x1*x2=1
  \draw[thick,blue!70!black,domain=0.52:2.35,samples=80,smooth]
    plot (\x,{1/\x});
  \node[blue!70!black,font=\small] at (2.0,0.62) {$x_1x_2=1$};
  % Iterates
  \coordinate (p0) at (2.0,1.0);
  \coordinate (p1) at (0.6,1.2);
  \coordinate (p2) at (1.007,0.852);
  \coordinate (p3) at (1.0,0.999);
  \coordinate (p4) at (1.0,1.0);
  % Arrows between iterates
  \draw[->,thick,red!70!black] (p0)--(p1);
  \draw[->,thick,red!70!black] (p1)--(p2);
  \draw[->,thick,red!70!black] (p2)--(p3);
  \draw[->,thick,red!70!black] (p3)--(p4);
  % Points
  \foreach \p/\lbl/\pos in {
    p0/$k=0$/above right,
    p1/$k=1$/left,
    p2/$k=2$/below right,
    p4/$x^*$/above right}{
    \fill[red!80!black] (\p) circle (1.6pt);
    \node[\pos,font=\tiny] at (\p) {\lbl};
  }
  % x* highlighted
  \fill[green!50!black] (p4) circle (2pt);
\end{tikzpicture}
\end{column}
\end{columns}

\end{frame}

\begin{frame}
\frametitle{Autres fonctions de pénalité}

Considérons le problème général avec contraintes
\begin{align*}
\min_{x \in \RR^n}\ & f(x) \\
\mbox{t.q. } & g_i(x) = 0,\ i \in \cE \\
& g_i(x) \leq 0,\ i \in \cI
\end{align*}

{\red Fonction de pénalité exacte}: $\Phi(x, \mu)$ est dite exacte s'il existe un $\mu^* > 0$ tel que si $0< \mu < \mu^*$, n'importe quel minimiseur local du problème avec contraintes est aussi minimiseur local de $\Phi(x, \mu)$.

\mbox{}

Note: la pénalité quadratique est inexacte.

\end{frame}

\begin{frame}
\frametitle{Autres fonctions de pénalité}

Exemples de pénalités exactes:
\begin{itemize}
\item 
fonction de pénalité $\ell_2$ si $\cI = \emptyset$:
$$
\Phi(x, \mu) = f(x) + \frac{1}{\mu} \| g(x) \|_2
$$
\item 
fonction de pénalité $\ell_1$:
notant $z^+ = \max\{z, 0\}$,
$$
\Phi(x, \mu) = f(x) + \frac{1}{\mu} \sum_{i \in \cE} | g_i(x) | + \frac{1}{\mu} \sum_{i \in \cI} [ g_i(x) ]^+
$$
\end{itemize}


Extension de la pénalité quadratique:
$$
\Phi(x, \mu) = f(x) + \frac{1}{2\mu} \| g(x) \|^2_2
+ \frac{1}{2\mu} \sum_{i \in \cI} \left( [ g_i(x) ]^+ \right)^2
$$
Il s'agit d'une pénalité inexacte, et non nécessairement douce (en particulier, elle peut ne pas être différentiable en tout point).

\end{frame}

\end{document}